\section{Reproducibility}
\begin{itemize}
  \item \textbf{Environment}: \texttt{uv} (\texttt{pyproject.toml} and \texttt{uv.lock}), Python 3.12;
        formatting and linting with \texttt{black} and \texttt{ruff}, enforced by a pre-commit hook.
  \item \textbf{Configuration}: every setting (paths, MLflow, data source, cleaning thresholds, features,
        split, hyper-parameters) is in \nolinkurl{.env}. It is loaded into a typed pydantic settings model with
        nested sections and validation, built once by a cached \texttt{get\_settings()}
        (\nolinkurl{src/weather/utils/config.py}). Environment variables override the file.
  \item \textbf{Pipeline}: \texttt{uv run dvc repro} rebuilds every artefact. Every stage depends on
        \nolinkurl{.env}, and \texttt{dvc metrics diff} compares experiments.
  \item \textbf{Data versioning}: raw and processed data and model binaries are DVC-tracked and pushed to a
        DVC remote. Figures, tables and metrics are committed to git.
  \item \textbf{Experiment tracking}: \texttt{uv run mlflow ui --backend-store-uri sqlite:///mlflow.db}
        shows one run per stage, nested runs per model, artefacts, and registered models.
  \item \textbf{Tests}: \texttt{uv run pytest}. The suite covers cleaning rules, corrupted-batch detection,
        alignment of features and targets, strict time ordering of the split (no leakage), and the metrics.
\end{itemize}
